\section{Introduction}
\label{sec:introduction}

Evaluation protocols carry claims. A held-out-form protocol for symbolic expressions (withhold
one equivalent form of each equation from training, then ask a learned encoder to identify it)
is understood to carry the claim that the encoder has recognised algebraic \emph{structure} rather
than surface tokens. The protocol does not establish that claim on its own. Something has to rule
out the cheaper explanations, and in this literature that step is usually skipped.

We ran the cheapest possible check on the canonical case. On a 10-equation held-out-form suite
built from AI~Feynman \citep{udrescu2020ai}, a trained Tree-LSTM identifies held-out forms at
$0.972 \pm 0.060$ ($n{=}25$ seeds, chance $0.10$), against an untrained encoder's
$0.724 \pm 0.184$. Read naively, that is structural identification. It cannot be. On this suite every equation has a
variable set no other equation shares (\S\ref{sec:survey}), so a bag over variable tokens alone
(nearest-centroid, \emph{no training and no parameters}) is a complete lookup solution \emph{by
construction}, and duly scores $1.000$ (sd $0.000$), above the trained encoder at every seed. A
full-token bag lands nearby but unstably: $0.992 \pm 0.027$ in this run, $0.932 \pm 0.075$ in an
independent replication of the same protocol, a spread that straddles $0.972$ and is why we anchor
on the deterministic row (\S\ref{sec:reporting}). Under this construction the benchmark is solvable
by lookup: whatever the encoder learned, this protocol cannot see it, because a method that provably
has no notion of tree structure already saturates the metric.

\textbf{Ten equations is a small suite, and this is a constructed counterexample, not evidence of
prevalence}: showing that a failure mode exists and can be total needs one case, not a large
sample. Prevalence is assessed separately by the 23-corpus audit (\S\ref{sec:survey}) and behaviour
at scale by a published leaderboard (\S\ref{sec:leaderboard}). With $K{=}10$ the unit of inference
also matters from the outset: per-equation leakage ranges from $-0.05$ to $1.12$, so the $\pm$
figures above are seed-level and understate equation-level uncertainty by roughly $6\times$; we
bootstrap over equations throughout (\S\ref{sec:reporting}).

This gives the paper its thesis. \emph{Held-out-form accuracy can be explained by progressively
stronger non-learned cues; we give a training-free audit and a controlled-probe protocol for
determining which level of explanation a benchmark actually rules out.} The failure mode throughout
is not a model that underperforms but a measurement whose connection to the claim it supports was
never checked.

\paragraph{A tier decomposition of what a benchmark can be solved by.} ``Structure'' is not one
thing. We separate variable identity (tier~1), operator and arity inventory (tier~2), order and
local tree shape (tier~3), and genuine composition (tier~4), and argue that identification claims
must name their tier (\S\ref{sec:tiers}). The distinction is not cosmetic: a benchmark can be fully
immune at tier~1 and still be separable at tier~2, which is the situation almost everywhere we
looked, and clearing tier~2 still leaves tier~3, where surface statistics of a tree substitute
for composing it.

\paragraph{Testing for what the lower tiers cannot explain.} We measure tier~3 with two
zero-parameter baselines (token n-grams; parent--child and subtree-size statistics), then hold both
inventory \emph{and} local shape fixed with a \emph{twin} probe: each held-out form paired with a
rearrangement of itself carrying identical token and operator multisets, so those baselines are
pinned at exactly chance by construction (\S\ref{sec:tier3}). This tests sensitivity to
compositional distinctions not captured by the measured tier-3 statistics; it does not isolate
composition, and \S\ref{sec:tier3} says why no such probe could. The trained encoder is near
ceiling there, but an \emph{untrained} encoder reaches $0.788$, which is why we call these
baselines \emph{non-learned} rather than ``structure-free'': a randomly initialised recursive
encoder computes over structure without having learned any.

\paragraph{A training-free screen, and a targeted audit of 23 corpora.} Tier separability is computable
from the corpus alone, with no model and no training run (\S\ref{sec:auditor}). We ship it as
\texttt{audit-symbolic-benchmark} and apply it to 23 published corpora across four families
(\S\ref{sec:survey}). The result narrows the cautionary claim considerably: the tier-1 flaw is a
property of a per-symbol \emph{tokenization convention} applied to physics corpora, not of
equivalence-class benchmarks in general. The canonical published family, EQNET
\citep{allamanis2017learning} and successors \citep{gangwar2023semantic, zheng2025egen},
defines classes over a shared variable pool and is immune at tier~1 ($\leq 3.2\%$ unique
variable sets across all 15 corpora). Tier~2 is where the exposure is.

\paragraph{A floor under a live leaderboard, and a retraction.} The strongest test of an audit is
to point it at numbers people are currently standing behind. Three published systems report
$\mathrm{score}_5$ on the \textsc{UnseenEqClass} split of the 14 SemVec corpora. The metric needs
only a distance, so it is computable for a token bag exactly as for a learned embedding. Computed
under those authors' own procedure, the audit yields a baseline floor under the published protocol: no non-learned
representation reaches EqNet or the headline SemEmb system on any of the 14 corpora
(\S\ref{sec:leaderboard}). We report this having first published a stronger claim internally and
retracted it: our own protocol, internally consistent and satisfying our own checklist, made
the same floor appear to beat a published ablation on 13 of 14 corpora instead of 5. Every
difference from the original procedure had pushed our numbers up. \emph{A skyline is only a
skyline for the evaluation it was computed under.}

\paragraph{What the audit costs the auditor.} We applied the standard to ourselves and it was not
free: supplying a baseline we had not previously measured repeatedly reversed conclusions we had
drawn, and two purpose-built ``leakage-proof'' suites turned out to be solvable at tier~2
(\S\ref{sec:reporting}). The unit of inference also matters: per-equation leakage ranges from
$-0.05$ to $1.12$, so a seed-level interval understates equation-level uncertainty by roughly
$6\times$.

We close with the eight-step checklist these findings imply (\S\ref{sec:checklist}).

\paragraph{Related work.} None of the machinery is new in kind. \citet{geirhos2020shortcut} frame
\emph{shortcut learning} as a pervasive failure mode; in NLI, hypothesis-only baselines matched full
models and exposed annotation artifacts rather than reasoning
\citep{gururangan2018annotation, poliak2018hypothesis}. Our zero-parameter bag is the symbolic
analogue, and on the physics suite it does not merely match the trained encoder but solves the
benchmark by construction. \citet{hewitt2019control} established with \emph{control tasks} that
probing accuracy is interpretable only against a randomised baseline, which our untrained-encoder
skyline supplies. Renaming-invariance has an analogue in code: Type-2 clones are by definition
identical up to identifier names, and clone benchmarks curate them as a distinct category
\citep{svajlenko2014towards}. SCAN \citep{lake2018generalization} and COGS \citep{kim2020cogs}
showed that how a held-out set is \emph{constructed} decides observed difficulty, and the
contamination literature \citep{jacovi2023stop, zhou2023don} recommends shortcut-breaking held-out
sets, which our checklist operationalises as a tool. What is new is the tier decomposition, the
training-free screen behind it, and the audit of which families are exposed at which tier.
